% Options for packages loaded elsewhere
\PassOptionsToPackage{unicode}{hyperref}
\PassOptionsToPackage{hyphens}{url}
\documentclass[
]{article}
\usepackage{xcolor}
\usepackage{amsmath,amssymb}
\setcounter{secnumdepth}{-\maxdimen} % remove section numbering
\usepackage{iftex}
\ifPDFTeX
  \usepackage[T1]{fontenc}
  \usepackage[utf8]{inputenc}
  \usepackage{textcomp} % provide euro and other symbols
\else % if luatex or xetex
  \usepackage{unicode-math} % this also loads fontspec
  \defaultfontfeatures{Scale=MatchLowercase}
  \defaultfontfeatures[\rmfamily]{Ligatures=TeX,Scale=1}
\fi
\usepackage{lmodern}
\ifPDFTeX\else
  % xetex/luatex font selection
\fi
% Use upquote if available, for straight quotes in verbatim environments
\IfFileExists{upquote.sty}{\usepackage{upquote}}{}
\IfFileExists{microtype.sty}{% use microtype if available
  \usepackage[]{microtype}
  \UseMicrotypeSet[protrusion]{basicmath} % disable protrusion for tt fonts
}{}
\makeatletter
\@ifundefined{KOMAClassName}{% if non-KOMA class
  \IfFileExists{parskip.sty}{%
    \usepackage{parskip}
  }{% else
    \setlength{\parindent}{0pt}
    \setlength{\parskip}{6pt plus 2pt minus 1pt}}
}{% if KOMA class
  \KOMAoptions{parskip=half}}
\makeatother
\setlength{\emergencystretch}{3em} % prevent overfull lines
\providecommand{\tightlist}{%
  \setlength{\itemsep}{0pt}\setlength{\parskip}{0pt}}
\usepackage{bookmark}
\IfFileExists{xurl.sty}{\usepackage{xurl}}{} % add URL line breaks if available
\urlstyle{same}
\hypersetup{
  pdftitle={Julia Programming},
  hidelinks,
  pdfcreator={LaTeX via pandoc}}

\title{Julia Programming}
\author{}
\date{}

\begin{document}
\maketitle

\subsection{Description}\label{description}

Julia is a high-level, high-performance programming language designed
for technical computing. It combines the ease of use of \texttt{Python}
with the speed of C, making it ideal for scientific computing, data
science, and \texttt{machine\ learning}. This course covers Julia from
its fundamentals through advanced topics including its powerful type
system, multiple dispatch, metaprogramming, parallel and \texttt{GPU}
computing, and the rich ecosystem of packages for data analysis and
machine learning. Students will learn to write performant, idiomatic
Julia code and leverage the language's unique strengths. The course
includes hands on exercises.

\subsection{Duration}\label{duration}

32 hours / 4 days

\subsection{Intended Audience}\label{intended-audience}

\begin{itemize}
\tightlist
\item
  Developers looking to adopt a high-performance language for technical
  computing.
\item
  Data scientists seeking a faster alternative to \texttt{Python} or R.
\item
  Researchers and engineers working on numerical simulation,
  optimization, or \texttt{machine\ learning}.
\end{itemize}

\subsection{Prerequisites}\label{prerequisites}

\begin{itemize}
\tightlist
\item
  Programming experience in at least one language such as
  \texttt{Python}, \texttt{MATLAB}, R, or C.
\item
  Basic understanding of linear algebra and statistics is helpful.
\item
  Familiarity with command-line interfaces.
\end{itemize}

\subsection{Required Knowledge}\label{required-knowledge}

\begin{itemize}
\tightlist
\item
  \texttt{Python} Programming (or equivalent experience)
\end{itemize}

\subsection{Objectives}\label{objectives}

\begin{itemize}
\tightlist
\item
  Write idiomatic and performant Julia code
\item
  Leverage multiple dispatch and the type system effectively
\item
  Use metaprogramming to generate efficient code
\item
  Perform parallel and \texttt{GPU} computing in Julia
\item
  Work with data using DataFrames.jl and visualization libraries
\item
  Build \texttt{machine\ learning} models with \texttt{Flux}.jl
\item
  Create and publish Julia packages
\end{itemize}

\subsection{Outline}\label{outline}

\begin{itemize}
\tightlist
\item
  Getting Started with Julia

  \begin{itemize}
  \tightlist
  \item
    What is Julia and why use it?
  \item
    Installation and environment setup
  \item
    The Julia REPL and workflow
  \item
    Variables, expressions, and basic types
  \item
    Strings and string interpolation
  \item
    Control flow: if/else, for, while
  \item
    First program and basic \texttt{I/O}
  \end{itemize}
\item
  Functions and Multiple Dispatch

  \begin{itemize}
  \tightlist
  \item
    Defining functions (standard and short form)
  \item
    Anonymous functions and closures
  \item
    Multiple dispatch fundamentals
  \item
    Method definitions and dispatch rules
  \item
    Variadic and keyword arguments
  \item
    Function composition and piping
  \end{itemize}
\item
  Type System

  \begin{itemize}
  \tightlist
  \item
    \texttt{Abstract-types} and type hierarchy
  \item
    Concrete types and primitive types
  \item
    Composite types (structs: mutable and immutable)
  \item
    Parametric types
  \item
    Union types
  \item
    Type annotations and conversions
  \item
    Designing with the type system
  \end{itemize}
\item
  Arrays and Linear Algebra

  \begin{itemize}
  \tightlist
  \item
    \texttt{Array} creation and initialization
  \item
    Indexing, slicing, and views
  \item
    Broadcasting
  \item
    Comprehensions and generators
  \item
    Matrix operations and linear algebra
  \item
    Sparse arrays
  \item
    The LinearAlgebra standard library
  \end{itemize}
\item
  Metaprogramming

  \begin{itemize}
  \tightlist
  \item
    Expressions and symbols
  \item
    Quoting and interpolation
  \item
    Macros: definition and usage
  \item
    Generated functions
  \item
    Code introspection (@code\_lowered, @code\_native)
  \item
    Common macros and when to use metaprogramming
  \end{itemize}
\item
  Modules and Packages

  \begin{itemize}
  \tightlist
  \item
    Modules: definition, export, and import
  \item
    The Pkg package manager
  \item
    Project environments and Project.\texttt{toml}
  \item
    Finding and using packages
  \item
    Creating your own package
  \item
    Package versioning and compatibility
  \end{itemize}
\item
  \texttt{I/O} and Data Formats

  \begin{itemize}
  \tightlist
  \item
    \texttt{File} \texttt{I/O} (text and binary)
  \item
    \texttt{CSV}.jl for tabular data
  \item
    \texttt{JSON}.jl and \texttt{YAML}.jl
  \item
    JLD2.jl for Julia-native serialization
  \item
    Working with databases
  \item
    \texttt{HTTP} and web requests
  \end{itemize}
\item
  Performance Tips and Optimization

  \begin{itemize}
  \tightlist
  \item
    Type stability and its importance
  \item
    Benchmarking with BenchmarkTools.jl
  \item
    Profiling and identifying bottlenecks
  \item
    Memory allocation and avoiding allocations
  \item
    Performance annotations (@inbounds, @simd, @fastmath)
  \item
    Writing C-speed Julia code
  \end{itemize}
\item
  Parallel and Distributed Computing

  \begin{itemize}
  \tightlist
  \item
    Multithreading with Threads.@threads
  \item
    Thread safety and atomic operations
  \item
    Task-based concurrency
  \item
    Distributed computing with Distributed.jl
  \item
    pmap and distributed arrays
  \item
    Parallel patterns and best practices
  \end{itemize}
\item
  \texttt{GPU} Computing

  \begin{itemize}
  \tightlist
  \item
    Introduction to \texttt{GPU} programming
  \item
    CUDA.jl for NVIDIA GPUs
  \item
    Writing \texttt{GPU} kernels
  \item
    \texttt{Array} operations on the \texttt{GPU}
  \item
    Performance considerations for \texttt{GPU} code
  \end{itemize}
\item
  Data Science with Julia

  \begin{itemize}
  \tightlist
  \item
    DataFrames.jl: creating and manipulating data frames
  \item
    Data transformation and aggregation
  \item
    Split-apply-combine operations
  \item
    Joins and reshaping
  \item
    Missing data handling
  \end{itemize}
\item
  Visualization

  \begin{itemize}
  \tightlist
  \item
    Plots.jl: basic and advanced plotting
  \item
    Plot backends and customization
  \item
    Makie.jl for interactive and publication-quality graphics
  \item
    Statistical visualizations
  \end{itemize}
\item
  \texttt{Machine\ Learning} with \texttt{Flux}.jl

  \begin{itemize}
  \tightlist
  \item
    Introduction to \texttt{Flux}.jl
  \item
    Building neural network layers
  \item
    Training loops and loss functions
  \item
    Optimizers and gradient descent
  \item
    Model evaluation and saving
  \item
    Common architectures (dense, convolutional, recurrent)
  \end{itemize}
\item
  Interoperability

  \begin{itemize}
  \tightlist
  \item
    Calling C and \texttt{Fortran} from Julia
  \item
    PyCall.jl: calling \texttt{Python} from Julia
  \item
    RCall.jl: calling R from Julia
  \item
    Wrapping external libraries
  \end{itemize}
\item
  Testing and Package Development

  \begin{itemize}
  \tightlist
  \item
    Unit testing with Test standard library
  \item
    Test sets and test organization
  \item
    Continuous integration for Julia packages
  \item
    Documentation with Documenter.jl
  \item
    Package development workflow
  \item
    Publishing to the General registry
  \end{itemize}
\end{itemize}

\subsection{Copyright}\label{copyright}

Mark Veltzer
\href{mailto:mark.veltzer@gmail.com}{\nolinkurl{mark.veltzer@gmail.com}},
© 2026

\end{document}
